\section{HeapSort}

\begin{softnote}
	HeapSort è l'ultimo algoritmo di ordinamento per confronti che
	studiamo. Combina il meglio di InsertionSort (ordina \textbf{sul posto})
	e di MergeSort (complessità $O(n \lg n)$), ma introduce una tecnica
	nuova: si basa su una \textbf{struttura dati} dedicata, lo \emph{heap},
	che trova applicazione anche nelle code di priorità.
\end{softnote}

% ----------------------------------------------------------------
\subsection{Idea Generale}

HeapSort:
\begin{itemize}
	\item ordina \textbf{sul posto} come InsertionSort;
	\item ha complessità $O(n \lg n)$ nel caso peggiore, come MergeSort;
	\item nel caso medio è \textbf{meno veloce} di QuickSort, perché le
	costanti nascoste in $O(n \lg n)$ sono più piccole in QuickSort
	che in HeapSort;
	\item è basato sulla struttura dati \textbf{heap}, usata anche in
	altri algoritmi importanti come le code di priorità.
\end{itemize}

% ----------------------------------------------------------------
\subsection{Struttura Dati Heap}

\begin{definizione}[Heap binario]
	Uno \textbf{heap} è un array che rappresenta un albero binario
	\emph{quasi completo}:
	\begin{itemize}
		\item ogni nodo dell'albero corrisponde a un elemento dell'array;
		\item tutti i livelli devono essere completamente riempiti tranne
		eventualmente l'ultimo, che può essere parziale ma sempre
		riempito da sinistra verso destra;
		\item se $A$ è uno heap: $A.\mathit{length}$ è il numero di elementi
		dell'array, $A.\mathit{heapSize}$ è il numero di nodi
		dell'albero rappresentati ($0 \le A.\mathit{heapSize} \le
		A.\mathit{length}$);
		\item la radice dell'albero è in $A[1]$.
	\end{itemize}
\end{definizione}

\begin{nota}
	\textbf{Albero binario completo vs.\ quasi completo:}
	\begin{itemize}
		\item \textbf{Completo:} ciascun nodo interno ha esattamente due
		figli; le foglie sono tutte allo stesso livello.
		\item \textbf{Quasi completo:} tutti i livelli riempiti tranne
		eventualmente l'ultimo, sempre riempito da sinistra.
		(Se si numerano i nodi per livelli, la numerazione non ha
		``buchi''.)
	\end{itemize}
\end{nota}

\subsubsection{Navigazione dell'albero tramite indici}

Dato il nodo in posizione $i$, genitore e figli si calcolano con:
\[
\textsc{Parent}(i) = \lfloor i/2 \rfloor, \qquad
\textsc{Left}(i)   = 2i,               \qquad
\textsc{Right}(i)  = 2i + 1.
\]

\begin{nota}
	Moltiplicare/dividere per 2 richiede solo di spostare di 1 bit a
	sinistra/destra: operazione estremamente efficiente a livello hardware.
\end{nota}

\subsubsection{Tipi di heap e proprietà}

\begin{definizione}[Max-heap e min-heap]
	\begin{itemize}
		\item \textbf{Max-heap:} $A[\textsc{Parent}(i)] \ge A[i]$ per ogni
		$i \ne 1$ (il genitore è sempre $\ge$ dei figli).
		\item \textbf{Min-heap:} $A[\textsc{Parent}(i)] \le A[i]$ per ogni
		$i \ne 1$ (il genitore è sempre $\le$ dei figli).
	\end{itemize}
	\textsc{HeapSort} usa i \textbf{max-heap}; i min-heap sono usati per
	realizzare le code di priorità.
\end{definizione}

\subsubsection{Altezza dello heap}

\begin{definizione}[Altezza di un nodo]
	L'\textbf{altezza} di un nodo è il numero di archi del cammino
	semplice più lungo che parte dal nodo e termina in una foglia
	(numero di livelli sottostanti al livello del nodo).
	L'\textbf{altezza dello heap} è l'altezza della radice.
\end{definizione}

\begin{teorema}[Altezza di un heap di $n$ nodi]
	Un albero binario quasi completo di $n$ nodi ha altezza
	$h = \lfloor \lg n \rfloor$.
\end{teorema}

\begin{proof}
	Vale $2^h \le n < 2^{h+1}$. Applicando $\lg$: $h \le \lg n < h+1$,
	quindi $h = \lfloor \lg n \rfloor$.
\end{proof}

\begin{esempio}[Altezza con $n=10$]
	$h = \lfloor \lg 10 \rfloor = \lfloor 3{,}32 \rfloor = 3$.
\end{esempio}

\begin{nota}
	Quasi tutte le operazioni fondamentali sugli heap sono eseguite in
	un tempo proporzionale all'altezza dell'albero, quindi $O(\lg n)$.
\end{nota}

\subsubsection{Operazioni fondamentali sullo heap}

Sono 6 le operazioni fondamentali:
\begin{enumerate}
	\item \textsc{BuildMaxHeap} — genera un max-heap da un array
	generico in tempo $O(n)$;
	\item \textsc{MaxHeapify} — ripristina la proprietà di max-heap
	quando a un elemento viene diminuito il valore; tempo
	$O(\lg n)$ (necessaria a \textsc{BuildMaxHeap});
	\item \textsc{MaxHeapInsert}, \textsc{HeapExtractMax},
	\textsc{HeapIncreaseKey} e \textsc{HeapMax} — consentono di
	usare uno heap per realizzare una coda di priorità; complessità
	$O(\lg n)$ tranne \textsc{HeapMax} che è $O(1)$.
\end{enumerate}

% ----------------------------------------------------------------
\subsection{MaxHeapify}

\begin{softnote}
	\textsc{MaxHeapify} è l'operazione di base: dato un nodo $i$ il cui
	valore è stato diminuito, ``fa scendere'' $A[i]$ fino alla posizione
	corretta ripristinando la proprietà di max-heap nel sottoalbero
	radicato in $i$.
\end{softnote}

\paragraph{Precondizione.}
I sottoalberi radicati in $\textsc{Left}(i)$ e $\textsc{Right}(i)$
sono già max-heap; la proprietà può essere violata solo nel nodo $i$.

\begin{algorithm}[H]
	\caption{\textsc{MaxHeapify}$(A, i)$}
	\KwIn{Array $A$, indice $i$; sottoalberi di $\textsc{Left}(i)$ e
		$\textsc{Right}(i)$ sono max-heap}
	\KwOut{$A$ con la proprietà di max-heap ristabilita nel sottoalbero
		radicato in $i$}
	$l \gets \textsc{Left}(i)$\;
	$\mathit{max} \gets (l \le A.\mathit{heapSize}\ \KwAnd\ A[l] > A[i])
	\mathbin{?}\ l : i$\tcp*{Figlio sx vs.\ radice}
	$r \gets \textsc{Right}(i)$\;
	\If{$r \le A.\mathit{heapSize}$ \KwAnd $A[r] > A[\mathit{max}]$}{
		$\mathit{max} \gets r$\;
	}
	\If{$\mathit{max} \ne i$}{
		scambia $A[i]$ con $A[\mathit{max}]$\;
		$\textsc{MaxHeapify}(A, \mathit{max})$\tcp*{Ricorsione sul figlio}
	}
\end{algorithm}

\subsubsection{Analisi di complessità di MaxHeapify}

Sia $n$ il numero di nodi del sottoalbero da sistemare.

\paragraph{Ricorrenza.}
\textsc{MaxHeapify} esegue $\Theta(1)$ operazioni più una sola chiamata
ricorsiva sul sottoalbero più grande:
\[
T(n) \le T\!\left(\tfrac{2}{3}n\right) + \Theta(1).
\]

\paragraph{Giustificazione del fattore $\frac{2}{3}$.}
Nel caso peggiore l'ultimo livello dell'albero è riempito esattamente
a metà: il sottoalbero sinistro ha un livello in più del destro. Se
$h$ è l'altezza della radice, il sottoalbero sinistro (completo) ha
$2^h - 1$ nodi, quello destro (completo) $2^{h-1} - 1$ nodi e
il totale è $n = 3 \cdot 2^{h-1} - 1$. La proporzione del sottoalbero
sinistro è:
\[
\frac{2^h - 1}{3 \cdot 2^{h-1} - 1}
\;\le\; \frac{2^h}{3 \cdot 2^{h-1}}
= \frac{2}{3}.
\]

\paragraph{Soluzione.}
Applicando il Teorema dell'Esperto (caso 2), la ricorrenza
$T(n) \le T\!\left(\alpha n\right) + \Theta(1)$ ha soluzione
$T(n) = O(\lg n)$ per qualsiasi $\alpha < 1$:
\[
T_{\textsc{MaxHeapify}}(n) = O(\lg n) = O(h).
\]

\begin{hardnote}
	\textbf{Caso peggiore concreto:} un elemento posto alla radice deve
	scendere fino a una foglia, eseguendo $\lfloor \lg n \rfloor$ chiamate
	ricorsive. Quindi la complessità è anche $\Omega(\lg n)$, cioè
	$\Theta(\lg n)$ nel caso peggiore.
\end{hardnote}

% ----------------------------------------------------------------
\subsection{BuildMaxHeap}

\begin{softnote}
	Per trasformare un array generico in un max-heap si osserva che
	tutti gli elementi in $A[\lfloor n/2 \rfloor + 1 : n]$ sono
	\textbf{foglie}, quindi già heap di un elemento. Basta applicare
	\textsc{MaxHeapify} a tutti i nodi interni, risalendo dall'ultimo
	verso la radice.
\end{softnote}

\begin{algorithm}[H]
	\caption{\textsc{BuildMaxHeap}$(A)$}
	\KwIn{Array $A$ di $n$ elementi}
	\KwOut{$A$ riordinato come max-heap}
	$A.\mathit{heapSize} \gets A.\mathit{length}$\;
	\For{$i \gets \lfloor A.\mathit{length}/2 \rfloor$ \KwTo $1$}{
		$\textsc{MaxHeapify}(A, i)$\;
	}
\end{algorithm}

\subsubsection{Correttezza di BuildMaxHeap}

\begin{definizione}[Invariante di BuildMaxHeap]
	All'inizio di ogni iterazione del ciclo \texttt{for}, ogni nodo di
	indice $i+1, \ldots, n$ è la radice di un max-heap.
\end{definizione}

\begin{proof}
	\begin{enumerate}
		\item \textbf{Inizializzazione:} prima del I ciclo, $i = \lfloor n/2
		\rfloor$. Tutti i nodi successivi sono foglie, dunque radici di
		max-heap di un solo elemento. \checkmark
		\item \textbf{Conservazione:} i figli del nodo $i$ hanno indice
		maggiore di $i$, quindi per l'invariante sono radici di
		max-heap. \textsc{MaxHeapify}$(A, i)$ ripristina la proprietà
		nel sottoalbero radicato in $i$, garantendo che all'iterazione
		$i-1$ i nodi $i, \ldots, n$ siano tutti radici di max-heap.
		\checkmark
		\item \textbf{Conclusione:} alla fine del ciclo $i = 0$, quindi il
		nodo 1 è la radice di un max-heap che comprende tutti i nodi da
		1 a $n$: $A$ è un max-heap. \checkmark
	\end{enumerate}
\end{proof}

\subsubsection{Complessità di BuildMaxHeap}

\begin{softnote}
	Una stima immediata è $O(n \lg n)$ (si chiama \textsc{MaxHeapify}
	$n/2$ volte, ciascuna $O(\lg n)$). Ma si può dimostrare un limite
	più stretto: $O(n)$.
\end{softnote}

Per ogni altezza $h'$ con $0 \le h' \le \lfloor \lg n \rfloor$,
il numero di nodi ad altezza $h'$ è al più
$n_{h'} = \lceil n / 2^{h'+1} \rceil$.

\begin{proof}[Dimostrazione per induzione]
	Le foglie ($h'=0$) sono $\lceil n/2 \rceil = \lceil n/2^{0+1} \rceil$.
	Se $n_{h'-1} \le \lceil n/2^{h'} \rceil$, allora
	$n_{h'} = \lceil n_{h'-1}/2 \rceil \le \lceil n/2^{h'+1} \rceil$.
\end{proof}

\textsc{MaxHeapify} usa tempo $O(h')$ per ciascun nodo ad altezza
$h'$. Il costo totale è:
\[
\sum_{h'=1}^{\lfloor \lg n \rfloor}
\left\lceil \frac{n}{2^{h'+1}} \right\rceil O(h')
\;\le\;
O\!\left(n \sum_{h'=0}^{\lfloor \lg n \rfloor} \frac{h'}{2^{h'}}
\right)
\;=\;
O\!\left(n \cdot \frac{1/2}{(1-1/2)^2}\right)
= O(n).
\]

\begin{nota}
	Si usa la somma della serie $\sum_{h'=0}^{\infty} h' x^{h'} =
	x/(1-x)^2$ con $x = 1/2$, ottenendo il valore finito 2.
\end{nota}

\begin{teorema}[Complessità di BuildMaxHeap]
	\[
	T_{\textsc{BuildMaxHeap}}(n) = O(n).
	\]
\end{teorema}

% ----------------------------------------------------------------
\subsection{Algoritmo HeapSort}

\begin{softnote}
	Una volta costruito il max-heap, l'elemento massimo si trova sempre
	in $A[1]$. L'idea è estrarlo iterativamente: si scambia la radice con
	l'ultimo elemento, si riduce l'\textit{heapSize} di 1 e si ripristina
	la proprietà con \textsc{MaxHeapify}. Al termine l'array è ordinato in
	modo crescente.
\end{softnote}

\begin{algorithm}[H]
	\caption{\textsc{HeapSort}$(A)$}
	\KwIn{Array $A$ di $n$ elementi}
	\KwOut{$A$ ordinato in loco}
	$\textsc{BuildMaxHeap}(A)$\;
	\For{$i \gets A.\mathit{length}$ \KwTo $2$}{
		scambia $A[1]$ con $A[i]$\;
		$A.\mathit{heapSize} \gets A.\mathit{heapSize} - 1$\;
		$\textsc{MaxHeapify}(A, 1)$\;
	}
\end{algorithm}

\subsubsection{Correttezza di HeapSort}

\begin{definizione}[Invariante di HeapSort]
	All'inizio di ogni iterazione del ciclo \texttt{for}, il sottoarray
	$A[1:i]$ è un max-heap contenente gli $i$ elementi più piccoli di
	$A[1:n]$, e il sottoarray $A[i+1:n]$ contiene gli $n-i$ elementi più
	grandi di $A[1:n]$, ordinati in modo crescente.
\end{definizione}

\subsubsection{Complessità di HeapSort}

\begin{itemize}
	\item \textsc{BuildMaxHeap}: $O(n)$.
	\item Ciclo \texttt{for}: $n-1$ chiamate a \textsc{MaxHeapify},
	ciascuna $O(\lg n)$.
\end{itemize}

\[
T_{\textsc{HeapSort}}(n) = O(n) + (n-1)\,O(\lg n) = O(n \lg n).
\]

\begin{center}
	\begin{tabular}{lll}
		\toprule
		\textbf{Caso} & \textbf{Input} & \textbf{Complessità} \\
		\midrule
		Peggiore & qualsiasi & $O(n \lg n)$ \\
		Medio    & qualsiasi & $O(n \lg n)$ \\
		Migliore & qualsiasi & $O(n \lg n)$ \\
		\bottomrule
	\end{tabular}
\end{center}

\begin{hardnote}
	\textbf{Da ricordare all'esame:} HeapSort è $O(n \lg n)$ in tutti i
	casi e ordina sul posto. Nel caso medio è più lento di QuickSort
	(costanti peggiori), ma non soffre del caso peggiore $\Theta(n^2)$ di
	QuickSort non randomizzato.
\end{hardnote}

% ----------------------------------------------------------------
\subsection{Code di Priorità}

\begin{softnote}
	Lo heap non serve solo per ordinare: è la struttura dati ideale per
	realizzare una \textbf{coda di priorità}, che consente di estrarre
	sempre l'elemento con priorità massima (o minima) in modo efficiente.
\end{softnote}

\begin{definizione}[Coda di max-priorità]
	Una \textbf{coda di max-priorità} è una struttura dati che gestisce
	un insieme $S$ di elementi, ciascuno associato a un valore detto
	\emph{chiave} o \emph{priorità}, in modo da estrarre l'elemento di
	priorità massima o modificare la priorità di un elemento nel modo più
	efficiente possibile. Analogamente si definisce la coda di
	\textbf{min-priorità}.
\end{definizione}

\subsubsection{Operazioni fondamentali}

Se $x$ è un elemento, $x.\mathit{key}$ è la sua chiave/priorità.
Le operazioni fondamentali sono:
\begin{itemize}
	\item \textbf{Insert}$(S, x)$: inserisce $x$ in $S$.
	\item \textbf{Max}$(S)$: restituisce l'elemento di $S$ con priorità
	massima.
	\item \textbf{ExtractMax}$(S)$: rimuove e restituisce l'elemento con
	priorità massima.
	\item \textbf{IncreaseKey}$(S, x, k)$: aumenta la chiave di $x$ a
	$k$ (con $k \ge x.\mathit{key}$).
\end{itemize}

\begin{nota}
	Se si vuole \textbf{diminuire} il valore di un elemento, è sufficiente
	farlo direttamente e poi invocare \textsc{MaxHeapify} per ripristinare
	la proprietà di max-heap.
\end{nota}

\subsubsection{Implementazione delle operazioni}

\begin{algorithm}[H]
	\caption{\textsc{Max}$(S)$}
	\KwIn{$S$ coda di max-priorità}
	\KwOut{Elemento con priorità massima}
	\Return $S[1]$\;
\end{algorithm}

Complessità: $O(1)$.

\begin{algorithm}[H]
	\caption{\textsc{ExtractMax}$(S)$}
	\KwIn{$S$ coda di max-priorità}
	\KwOut{Elemento con priorità massima, rimosso da $S$}
	\If{$S.\mathit{heapSize} < 1$}{\Return \texttt{"Coda vuota!"}\;}
	$\mathit{max} \gets S[1]$\;
	$S[1] \gets S[S.\mathit{heapSize}]$\;
	$S.\mathit{heapSize} \gets S.\mathit{heapSize} - 1$\;
	$\textsc{MaxHeapify}(S, 1)$\;
	\Return $\mathit{max}$\;
\end{algorithm}

Complessità: $O(\lg n)$.

\begin{algorithm}[H]
	\caption{\textsc{IncreaseKey}$(S, i, \mathit{key})$}
	\KwIn{$S$ coda di max-priorità, posizione $i$, nuovo valore
		$\mathit{key}$}
	\KwOut{Priorità $S[i].\mathit{key}$ aggiornata in $S$}
	\If{$\mathit{key} < S[i].\mathit{key}$}{
		\Return \texttt{"Errore: nuova priorità minore dell'attuale"}\;
	}
	$S[i].\mathit{key} \gets \mathit{key}$\;
	\While{$i > 1$ \KwAnd $S[\textsc{Parent}(i)].\mathit{key}
		< S[i].\mathit{key}$}{
		scambia $S[i]$ con $S[\textsc{Parent}(i)]$\;
		$i \gets \textsc{Parent}(i)$\;
	}
\end{algorithm}

\begin{softnote}
	Se la nuova chiave rompe la proprietà di max-heap, l'elemento viene
	scambiato con il genitore fino a trovare la posizione corretta
	(risalita verso la radice).
\end{softnote}

Complessità: $O(\lg n)$ (al più si sale dalla foglia più profonda
alla radice).

\begin{algorithm}[H]
	\caption{\textsc{Insert}$(S, x)$}
	\KwIn{$S$ coda di max-priorità, nuovo elemento $x$}
	\KwOut{$S$ aggiornato con $x$}
	$S.\mathit{heapSize} \gets S.\mathit{heapSize} + 1$\;
	$\mathit{key} \gets x.\mathit{key}$\;
	$x.\mathit{key} \gets -\infty$\tcp*{Inserisce in fondo con priorità minima}
	$S[S.\mathit{heapSize}] \gets x$\;
	$\textsc{IncreaseKey}(S,\, S.\mathit{heapSize},\, \mathit{key})$\tcp*{Risale alla posizione corretta}
\end{algorithm}

Complessità: $O(\lg n)$ (l'elemento potrebbe risalire dalla foglia
più profonda alla radice).

\begin{center}
	\begin{tabular}{ll}
		\toprule
		\textbf{Operazione} & \textbf{Complessità} \\
		\midrule
		\textsc{Max}           & $O(1)$      \\
		\textsc{ExtractMax}    & $O(\lg n)$  \\
		\textsc{IncreaseKey}   & $O(\lg n)$  \\
		\textsc{Insert}        & $O(\lg n)$  \\
		\textsc{BuildMaxHeap}  & $O(n)$      \\
		\textsc{MaxHeapify}    & $O(\lg n)$  \\
		\bottomrule
	\end{tabular}
\end{center}



% ============================================================
